\section{Conclusion}

We began by observing a counterintuitive finding: two architectures designed for the same purpose produce measurably different internal representations, yet this difference translates to performance advantages only for global statistics tasks, not for local anomaly detection. This partial confirmation provides the first empirical guidance for practitioners selecting weight-space architectures.

\subsection*{Summary}

Our main contributions are:

\begin{enumerate}
    \item \textbf{First quantitative measurement of inductive bias differences:} DWS CoV 1.44 vs NFT CoV 1.35

    \item \textbf{Empirical confirmation of task-dependent advantage:} NFT RMSE 82.9 vs DWS 94.5 (12.3\% improvement)

    \item \textbf{Methodology for architecture comparison:} Controlled 2$\times$3 factorial design
\end{enumerate}

\subsection*{Future Directions}

\begin{itemize}
    \item \textbf{Testing Untested Alternative Explanations:} Train NFT with extended epochs to test if inductive bias difference narrows
    \item \textbf{Verifying Unconfirmed Assumptions:} Analyze real TrojAI backdoored models for spatial weight structure
    \item \textbf{Extending Scope:} Apply to Transformer weight-spaces and challenging datasets
\end{itemize}

Our findings suggest that the right tool for the right job matters in weight-space learning: architecture selection should consider how inductive biases align with task characteristics.
